% ============================================================
% PRM393 - Section 5: Multimedia & Real-Time (Bilingual)
% ============================================================
\section{\bi{Trình phát đa phương tiện và tương tác trực tiếp}{OmniPlayer \& Real-Time Multimedia Delivery}}
\label{sec:streaming}

\bi{
Trung tâm trải nghiệm nghe nhìn của ứng dụng di động là thành phần
\texttt{OmniPlayer} (\texttt{lib/presentation/widgets/omni\_player.dart}).
Trình phát này được thiết kế theo mô hình đa tầng tiếp nhận (Multi-Engine
Playback Architecture), có khả năng chuyển đổi liền mạch giữa luồng trực tiếp
độ trễ thấp từ MediaMTX, video ghi hình theo yêu cầu (VOD) và luồng nhúng bên
ngoài.
}{
At the epicenter of the mobile user experience is \texttt{OmniPlayer}
(\texttt{lib/presentation/widgets/omni\_player.dart}). Architected around a
multi-engine playback paradigm, it dynamically transitions across low-latency
MediaMTX live broadcasts, on-demand video (VOD), and external embed streams.
}

\subsection{\bi{Kiến trúc phát trực tiếp qua MediaMTX và WebRTC/HLS}{MediaMTX Low-Latency Streaming Pipeline}}

\bi{
Đối với các kênh phát trực tiếp, OmniCast kết hợp giao thức ingest WHIP tại
phòng thu (Live Studio) với máy chủ chuyển tiếp \textbf{MediaMTX}. Trên ứng
dụng di động, \texttt{OmniPlayer} tích hợp gói \texttt{video\_player} chính
thức của Flutter, tự động phân giải địa chỉ IP máy chủ trong mạng cục bộ (LAN
Host Discovery) để khắc phục tình trạng thiết bị thật không thể truy cập
\texttt{localhost}:
}{
For live channels, OmniCast couples studio WHIP ingestion with the high-concurrency
\textbf{MediaMTX} streaming engine. On the mobile client, \texttt{OmniPlayer}
leverages Flutter's official \texttt{video\_player} package, dynamically
resolving LAN host IP addresses to allow physical test devices to seamlessly
stream feeds without hardcoded local loopback constraints:
}

\begin{lstlisting}[language=Dart, caption={\bi{Cơ chế cấu hình phân giải địa chỉ phát trực tiếp động.}{Dynamic stream endpoint resolution in OmniPlayer.}}]
String resolveStreamUrl(String rawUrl, String hostIp) {
  if (rawUrl.contains('localhost') || rawUrl.contains('127.0.0.1')) {
    return rawUrl
      .replaceAll('localhost', hostIp)
      .replaceAll('127.0.0.1', hostIp);
  }
  return rawUrl;
}
\end{lstlisting}

\bi{
Trường hợp luồng phát là liên kết nhúng từ YouTube hoặc các đài truyền hình
đối tác, \texttt{OmniPlayer} tự động chuyển sang chế độ nhúng qua
\texttt{webview\_flutter}, loại bỏ sự phụ thuộc vào các thư viện bên ngoài cồng
kềnh và tránh xung đột mã nguồn nhị phân gốc.
}{
When rendering partner broadcast embeds (such as YouTube feeds),
\texttt{OmniPlayer} seamlessly mounts a headless \texttt{webview\_flutter}
canvas, eliminating bloated native library footprints and preventing binary
dependency collisions.
}

\subsection{\bi{Lớp phủ tương tác thời gian thực hai chiều}{Bi-Directional Interactive Overlay}}

\bi{
Xung quanh khung hình video, ứng dụng cung cấp ba phân hệ tương tác thời gian
thực:
}{
Surrounding the video viewport, the mobile client embeds three real-time
interaction layers:
}

\begin{enumerate}[leftmargin=*,itemsep=2pt]
  \item \textbf{\bi{Bình luận trực tiếp (Live Chat):}{Live Messaging:}}
  \bi{Đồng bộ thông qua WebSocket namespace \texttt{/live}. Tin nhắn mới xuất
  hiện ngay lập tức theo dạng danh sách cuộn tự động với hiệu ứng làm mờ nền
  bán trong suốt (translucent blur), không che khuất nội dung phát
  sóng.}{Synced via WebSocket namespace \texttt{/live}. New messages arrive
  instantly in an auto-scrolling semi-transparent overlay that never obscures
  video content.}

  \item \textbf{\bi{Thanh cảm xúc nổi (Floating Reactions):}{Animated Floating Reactions:}}
  \bi{Người xem có thể thả các biểu tượng cảm xúc (tim, vỗ tay, bất ngờ).
  Các cảm xúc được phát tán tức thì tới toàn bộ người xem qua WebSocket và
  kích hoạt hiệu ứng bay lên hoạt họa mượt mà bằng CustomPainter.}{Viewers
  can dispatch emoji reactions (hearts, claps, fire). Reactions broadcast to
  all concurrent viewers via WebSockets, rendering upward floating animations
  via Flutter CustomPainters.}

  \item \textbf{\bi{Hộp thoại quyên góp trực tiếp (Live Donation Modal):}{Live Donation Sheet:}}
  \bi{Tích hợp trong \texttt{live\_donation\_modal.dart}, cho phép người xem
  chọn các gói ủng hộ nhanh (5.000đ, 20.000đ, 50.000đ...) kèm lời nhắn gửi
  đến người dẫn chương trình. Thao tác được xác thực bảo mật và hiển thị hiệu
  ứng vinh danh nổi bật trên luồng phát.}{Embedded in
  \texttt{live\_donation\_modal.dart}, enabling viewers to tip creators with
  preset amounts and custom cheers. Contributions trigger celebratory banner
  animations on the live screen.}
\end{enumerate}

\subsection{\bi{Khả năng tự phục hồi khi kết nối gián đoạn}{Resilient Connection Recovery}}

\bi{
Khi luồng phát sóng bị đứt quãng do người dùng di chuyển vào vùng sóng yếu,
\texttt{OmniPlayer} kích hoạt thuật toán \emph{Exponential Backoff with
Jitter}: tự động thử kết nối lại sau 1s, 2s, 4s, 8s (tối đa 30s) mà không làm
đóng băng ứng dụng hay đưa người xem ra khỏi màn hình hiện tại.
}{
When network degradation interrupts a live transmission, \texttt{OmniPlayer}
triggers an \emph{Exponential Backoff with Jitter} retry strategy: attempting
stream reconnections after 1s, 2s, 4s, 8s intervals (capped at 30s) without
freezing UI thread execution or prematurely kicking the user off the player
screen.
}
